\section{秦：统一来得很快，帝国来得更快}

战国末年，六国还在，地图上没有哪一国自愿承认自己将被抹去。可秦已经拥有一种不同的节奏：关中能提供相对稳定的腹地，县邑和军功把人口、土地与兵役接得更紧，长期战争又使它比对手更熟悉怎样把一次胜利变成下一次战争的准备。

\subsection{东出以前：边地的压力与宫廷的收束}

这条节奏并非从统一前夕突然出现。前五世纪末，魏军沿黄河逼迫关西，秦一面在河西调整车骑、步卒和城邑守备，\eventanchor{WQ-0416-QIN-HEXI-CAVALRY}{秦·河西整军}一面逐步形成以赏罚、军功重整兵员的献公传统，\eventanchor{WQ-0384-QIN-XIANGONG-MILITARY-LAW}{秦·献公军法整顿}。它们未能立刻扭转对魏的劣势，却把“能否由君主持续调动兵员”的问题留给孝公。商鞅在前359年前后以变法、徙木取信等故事争取授权，\eventanchor{WQ-0359-QIN-REFORM-DEBATE}{秦·变法取信}；这些著名对话带有后出政治寓言的色彩，但新法必须取得执行信誉、并承受旧贵族阻力的处境并不因此消失。秦惠文王称王以后，王命、相印和使节文书也开始以对等姿态进入列国往来，\eventanchor{WQ-0325-QIN-KINGSHIP-DOCUMENTS}{秦·称王文书}周王室的名分中心遂更难维持。

长期东进也有其逆流。前252年前后，秦军不得不逐步解除邯郸方向的攻势，\eventanchor{WQ-0252-QIN-HANDAN-WITHDRAWAL}{秦·邯郸撤军}，救赵联军、兵疲与补给使继续围攻的收益下降。长平所得并未丧失，但这次停顿提醒人们：秦军并非只要向东移动便可自动取胜，河东、上党与关中之间的粮道仍须修整。后来统一战争所显出的连续性，正是在这种受挫、补给与再集结的循环里长成的。

前250年的继承尤其显示国家机器并不只靠战场维持。孝文王即位三日即卒，\eventanchor{WQ-0250-QIN-XIAOWEN-THREE-DAYS}{秦·孝文王三日崩}华阳夫人被尊为太后，\eventanchor{WQ-0250-HUAYANG-QUEEN-MOTHER}{秦·华阳太后}庄襄王继位，\eventanchor{WQ-0250-ZHUANGXIANG}{秦·庄襄王即位}宫廷迅速接管文书和将领指挥，\eventanchor{WQ-0250-QIN-SUCCESSION-ADMIN}{秦·王位交接}。吕不韦以相国入居决策中心，\eventanchor{WQ-0250-LVBUWEI-PREMIER}{秦·吕不韦为相}又受封文信侯、食蓝田十二县，\eventanchor{WQ-0250-LVBUWEI-WENXIN-FIEF}{秦·文信侯封邑}使门客、财赋与相位暂时合在一个网络中。传世叙事把这一网络写得格外戏剧化，仍可看出短促王年没有造成全面失序；它为年幼的嬴政保留了一个可运行、却尚未完全归于王权的朝廷。

秦在关外的落脚点同时前移。蒙骜取成皋、巩并初置三川郡，\eventanchor{WQ-0249-QIN-SANCHUAN-COMMANDERY}{秦·置三川郡}把函谷关外的道路、河洛渡口纳入常设转运，\eventanchor{WQ-0249-SANCHUAN-TRANSPORT}{秦·三川转运}。庄襄王去世后，十三岁的王政即位，\eventanchor{WQ-0247-ZHUANGXIANG-DEATH}{秦·庄襄王去世}\eventanchor{WQ-0247-QINZHENG}{秦·王政即位}相国与宗室将领继续辅政。其后骊山陵园长期营建，\eventanchor{WQ-0246-QIN-MAUSOLEUM}{秦·骊山陵园}郑国渠也在泾水灌区动工并逐步发挥效用，\eventanchor{WQ-0246-ZHENGGUO-CANAL}{秦·郑国渠}工程征用劳力、土石和水工知识，\eventanchor{WQ-0246-ZHENGGUO-CANAL-LABOR}{秦·郑国渠劳役}不能只从后来“富国强兵”的结果倒推为毫无代价的成功。它们一面扩大供给，一面使关中家庭更深地进入国家工程。

东向的城链也在被一点点拆开。蒙骜前244年取畅、有诡等城，\eventanchor{WQ-0244-MENGAO-CHANG}{秦·蒙骜取畅}前242年又取魏二十城、置东郡，\eventanchor{WQ-0242-MENGAO}{秦·置东郡}\eventanchor{WQ-0242-MENGAO-TWENTY-CITIES}{秦·魏地二十城}。新郡必须清点城邑、道路与赋役，\eventanchor{WQ-0242-DONGJUN-COUNTY-ADMIN}{秦·东郡接管}才可能成为大梁方向的粮运与征兵支点；这不是一份今已保存的户籍册能逐项证实的过程。蒙骜约在前240年去世，\eventanchor{WQ-0240-MENGAO-DEATH}{秦·蒙骜身后}蒙氏的军功资源却仍为王政时期北方与东进的将领来源。

少年君主与摄政网络终究不能长久并立。嫪毐先受封长信侯，形成围绕太后的私门势力，\eventanchor{WQ-0239-LAOAI-CHANGXIN}{秦·长信侯网络}前238年王政在雍行冠礼、取得亲政资格，\eventanchor{WQ-0238-QIN-YONG-CORONATION}{秦·雍城冠礼}嫪毐随即以太后宫、兵符和禁卫为争夺焦点发动叛乱，\eventanchor{WQ-0238-LAOAI}{秦·嫪毐之乱}\eventanchor{WQ-0238-LAOAI-PALACE-GUARD}{秦·宫禁兵符之争}。私生子传说使这次冲突被写成继承秩序危机，\eventanchor{WQ-0238-LAOAI-CHILDREN}{秦·后宫继承危机}而诛三族、迁徙党羽等惩罚又把清理扩及整个政治网络，\eventanchor{WQ-0238-LAOAI-KIN-PUNISHMENT}{秦·惩治嫪毐党羽}；具体人数、刑罚和宫闱细节都不应当作司法统计。

翌年逐客令虽一度针对非秦籍游士，李斯上书后即被收回，\eventanchor{WQ-0237-LISI}{秦·收回逐客令}客卿继续进入文书、战略和郡县事务，\eventanchor{WQ-0237-LISI-GUEST-OFFICIALS}{秦·客卿留秦}显示秦既疑惧外来结党，又离不开其知识。王翦、桓齮等前236年攻赵，取邺、安阳并压迫漳河、太行之间的联系，\eventanchor{WQ-0236-QIN-ZHAO-HEYE}{秦·进取赵南}；吕不韦在前235年被迁蜀而死，\eventanchor{WQ-0235-LUBUWEI-EXILE}{秦·吕不韦放逐}其门客受到解散、监视或吸纳，\eventanchor{WQ-0235-LUBUWEI-RETAINERS}{秦·收束吕氏门客}封邑和文书关系亦被重新分配，\eventanchor{WQ-0235-QIN-LVBUWEI-AFTERMATH}{秦·重组相国网络}。至前231年韩国献南阳，\eventanchor{WQ-0231-NANYANG}{秦·韩国献南阳}秦已在新郑前方取得一处可以试验接管、也可以继续推进的前沿。

\begin{sourcenote}[东出前的材料边界]
\sourcebook{史记·秦本纪}、\sourcebook{商君列传}、\sourcebook{吕不韦列传}、\sourcebook{秦始皇本纪}和各国世家提供上述战争、继承与人事变动的编年骨架；《战国策》与后出列传所保存的徙木、宫闱、上书和廷议言辞，则更适于说明后人怎样理解权力冲突，不能逐句还原现场。郑国渠、兵器、城址与关中灌区的考古研究能够检验工程和军事条件，却不能替三川、东郡或吕氏门客重建完整的账簿。这里所能确认的是秦逐渐取得转运、动员和接管的能力；每一项措施在何处以何种强度落地，仍有材料空白。
\end{sourcenote}

\subsection{六个国家，六种约束}

\eventanchor{QIN-0230-UNIFICATION}{秦·统一六国}

统一战争先落在最难退让的中原通道上。前230年，内史腾入韩，韩王安被俘，韩国遂为秦所并，\eventanchor{QIN-0230-HAN}{秦·灭韩}。新郑一带的城邑、人口和道路原是韩赖以周旋的本钱，却也使它始终贴近秦东出的方向；从伊阙受挫到上党与本土被隔断的压力，早已削窄了这个小国的选择。秦在此处得到的并非一座孤城，而是进入中原腹地的一段行政和交通网络。史书没有留下足以逐日复原的攻韩战报，因而不能把韩亡写成一次城下决斗；能够把握的是，长期的地缘挤压和秦国持续东进构成结构性原因，前230年的进军则使这一脆弱位置终于失去回旋余地。

韩王被俘后的难题是让道路继续通行。秦接收新郑的官吏、户籍、仓廪并处理旧王族象征，\eventanchor{WQ-0230-HAN-POST-CONQUEST}{秦·接收韩地}继而以新郑为节点设置颍川郡，\eventanchor{WQ-0230-YINGCHUAN-COMMANDERY}{秦·置颍川郡}使韩旧都成为赵、魏方向的后勤前沿。接管手续的细部没有成套文书留存，不能想象为各县同日换印；但灭韩战果若不能转成征税、传令和储粮的行政单位，秦的主力便仍会被堵在函谷关外。

韩国的困境也说明，为何秦最先取得的不是一块可以安静休整的战利品。地处通道固然便于接收城邑、粮食和道路，却使新得区域马上成为继续向东行军的前沿；守备、转运和重新安排官吏必须与下一场战役并行。对韩廷而言，可退让的纵深有限，能够换取援助或重组的时间也较少；对秦廷而言，若不能把攻取转化为稳定的补给和通行，东出的门户仍会受阻。这些是由地理和战局提出的约束，不应补写成我们没有材料可以证明的秘密谋划。

秦没有因此转向同其余诸国一决胜负，而是继续把对手拆开。赵在长平后仍能守住邯郸并维持朝廷三十余年，但精锐、财政与外交余地已难恢复；前228年，秦军破赵都，俘赵王迁，\eventanchor{QIN-0228-ZHAO}{秦·破赵都}。这场覆亡的直接后果是赵的主体政权消失，却不等于赵地即时尽平：公子嘉在代地维持残余中心，直到前222年才被秦军消灭。把长平直接说成赵亡，或把前228年的攻取说成毫无残留的全境占领，都会抹去这段间隔；秦要付出的代价，仍是向北方山地和代地延伸补给、逐一接管城邑。

对秦而言，邯郸的陷落解决的是平原上的政治中心，却没有替北方完成接管。赵的残余能在代地延续，正因都城的失守与边地的服从不是同一件事：山地、较长的运输距离和仍可聚集的人群，会把一场破都之战拖成数年的边地经营。故而前228年至前222年的间隔不是史书遗漏的一块空白，而是统一战争的成本。秦军在这里不能只计算斩获，还得留下能够守关隘、接收粮食并传递文书的人手；否则新开的北方通道会反过来成为后方的漏洞。传世纪年能够确定都城失守和代地最终覆亡的大致次序，不能替我们列出每一处关隘的驻兵、每批粮食的起运日期。

前227年荆轲刺秦失败后，秦王加强宫禁、命王翦等东向伐燕，\eventanchor{WQ-0227-QIN-WANG-RESPONSE}{秦·刺秦后的伐燕}把一场外交绝境转成北方全面战争。宫门内的动作被《刺客列传》写得极富戏剧性，难以逐帧采信；较稳固的是刺杀并未挽救燕，反而使秦能把军事压力进一步推向易水以北。

赵都失守后，王贲转而面对大梁。魏的早期强盛正来自居中的河流、道路和城邑，到了前225年，这一地理既使大梁成为坚城，也使周边水网成为可被利用的弱点。传世叙事称王贲引河、沟之水灌城，魏王假降，\eventanchor{QIN-0225-WEI}{秦·灌大梁灭魏}。水攻不是单凭一道命令即可完成的奇计：它以围城、堤防、水道与持续守军为使能条件，也把破坏转嫁给城内外居民和中原耕地。《史记》所记三月、引入何水以及城坏的具体过程，在不同叙事中不宜当作已经测绘的工程记录；但大梁以水攻失守、魏国退出列国的主线，显示秦能够把中原的交通环境转为围城手段。

这场战役特别显出资源与地理如何互相限制。围城者要在水势改变之前维持外圈，防备援军，也要让自己的粮道和营地不被同一片水网切断；守城者虽可借城垣与水道拒敌，却会在水位、粮食、疫病和城外耕地受损时被迫重新计算。因而“灌大梁”可以作为秦工程化攻城能力的例子，不能被浪漫化为一招制胜的传奇。它留下的不只是魏王投降的结果，还有一片需要重新清点、修复并纳入秦县廷征发的中原空间。关于决堤位置、持续时日以及对土地的破坏规模，现有传世材料不足以给出精确结论。

王贲同年还取楚北十余城，\eventanchor{WQ-0225-WANGBEN-NORTHERN-CHU}{秦·取楚北城}说明攻魏和南线施压并非彼此无关；但城名、路线和每一支军队的转调并无完整清单。前224年，王翦在楚境久驻、积粮，\eventanchor{WQ-0224-QIN-CHU-MOBILIZATION}{秦·伐楚积粮}又以“请六十万”表明远征须改用更大兵力与更长准备，\eventanchor{WQ-0224-WANGJIAN-SIXHUNDREDTHOUSAND}{秦·王翦大军伐楚}。“六十万”很可能是后出叙事用以强调规模的整齐数字，仍提示这不是一场可由速胜方案承担的南征。

南面则不能沿用同样的节奏。楚的纵深、人口与江淮之间的道路，使它不像韩、魏那样能因一座都城受压便立刻失去行动能力；秦前224年\eventref{QIN-0224-CHU}{战国·李信伐楚受挫}，正提醒人们不可把楚写成等待收割的残局。次年王翦、蒙武再进，楚王负刍被获，\eventanchor{QIN-0223-CHU}{秦·灭楚}。这次转折的近期条件，是秦愿意为更审慎的南征重新集中兵力、粮运和将领；楚方面则须在广阔疆域中维持调度，并承受此前长期秦楚战争所累积的压力。《史记·白起王翦列传》著名的王翦请“六十万”之说，适合提示伐楚所需动员的规模和后勤问题，却不能把这个整齐数字当作军籍总数，更不能将列传中的君臣言辞还原为军前实录。

李信的受挫与王翦随后取胜，不宜被压缩为两位将领性格的对照。楚地的纵深意味着军队离关中越远，粮食、替补、讯息和伤病的处理越依赖一条不能中断的转运链；楚军则可以在尚未失去全部腹地时保存行动空间。秦改由王翦主军，是对这种约束的调整：它承认南征须以更多可调动的人力、较长准备和较稳的补给来换取胜算。所谓“六十万”正因为过于整齐，更该看作后出叙事对大规模动员的强调；前224年受挫、前223年楚亡的年代主线较稳，军队确数、战场展开和个别谏言则仍须悬置。

负刍被俘使楚主体政权覆亡，\eventanchor{QIN-0223-CHU-FUCHU}{秦·俘楚王负刍}却没有使淮南、江汉立刻安静。传世材料称昌平君战死，\eventanchor{WQ-0223-CHANGPINGJUN-DEATH}{秦·昌平君战死}其身份、地点与时间尚有争论；较可把握的是秦须继续接收城邑、压制残余抵抗，\eventanchor{WQ-0223-CHU-OCCUPATION}{秦·接管楚地}不能把俘王等同于占领结束。

楚亡后，秦得以收束北方的残余。燕在\eventref{QIN-0226-YAN}{秦·破蓟}后，燕王喜退守辽东；辽东的距离延缓了结局，却无法替失去都城、粮道和盟友的政权重建军力。前222年，王贲至辽东，燕王喜被俘，\eventref{QIN-0222-YAN}{秦·灭燕}。同年秦又灭代地残赵，说明燕亡并不是一支军队在地图边缘顺手完成的附带动作，而是秦在赵、燕两处残余之间继续分配行军与接管力量的结果。《秦始皇本纪》与《燕召公世家》保存了这一结果和大致次序；王贲军的逐日路线、各城何时开降，以及燕廷内部如何决定，材料都不足以写成无缝的战地纪录。

燕的末路把“距离”显示为双重因素。辽东远离秦的核心腹地，确实能延长行军和补给，也让残余朝廷暂时躲开中原决战；但距离同样把它同旧都、旧粮道和可能的援手隔开。秦在前226年破蓟、前222年处理辽东残部之间，并没有得到一条自动闭合的北方边界，反而必须同时顾及代地与燕的两个尾端。后人若把辽东写成燕王一退即可安然据守，便把边地的运输、城邑接管和兵源补充都从地图上抹去了。

王贲在前222年并灭辽东燕与代地赵的残部，\eventanchor{WQ-0222-NORTH-EAST-CLEARANCE}{秦·清除燕赵残部}才使北方不再留下可借山海和宗室重组的复国中心；两战先后与行军路线仍有异说。它释放的不是一支无需补给的“闲军”，而是秦在齐地投降以前得以重新调配的守备、道路和行政压力。

至此尚存齐。临淄的财富、海滨和半岛腹地仍可供给一国，但齐在秦赵、秦魏、秦楚等战争中长期未能形成可靠的共同防线，等到邻国相继覆亡，东方的资源已难转化为外援。前221年，王贲军自北而南入齐，齐王建降，\eventanchor{QIN-0221-QI}{秦·齐降秦}，六国至此尽亡。后出的叙事常把齐的孤立归结为君王后、后胜或若干宾客的一两次劝说，也喜欢把齐军未战的形象写得格外整齐；这些说法可以揭示齐廷在外交判断上的困境，不能替代地方守备、道路控制和降附过程的档案。较稳妥的解释是，齐的失败是长期孤立这一结构性条件、秦军由北面进入的直接压力，以及其余五国已不能牵制秦的共同结果。六次覆亡因而不是同一种胜利的重复：秦每次都利用不同国家的地理、兵力和联盟困局，又用前一次接收的道路、粮食与人力缩短下一次战争。

齐的降服并不证明半岛富庶在军事上毫无价值，而是说明资源必须经由盟友、道路、守备和决策才能成为可用的力量。当北面的旧赵、燕方向已由秦控制，其他五国又不复存在，齐所能调用的资源已难以转化成解除包围或牵制敌军的外部条件。秦军能够由北面推进，依靠的也不是一纸“齐降”的诏令，而是此前数次战争打开的通道和对中原节点的控制。关于齐军是否全然未战、哪些城邑先降、齐廷中谁负主要责任，传世叙事远比现存地方材料丰富；宜将其当作汉人解释亡国的语言，而不当作可逐项核实的作战日志。

\begin{causalchain}[统一战争的递进]
韩的覆亡让秦越过最狭窄的中原门户；赵、魏的失败清除了北方平原和水网中的两个强邻；伐楚受挫后改以更大后勤投入南征，表明秦也必须向地理让步；燕、代与齐的最终处置，则利用了前面战役已经改变的道路、盟友和资源分布。这个递进是使能条件，不是说每一次灭国都由上一次机械推出。各国的朝廷、守军和地方社会仍有自己的选择，现存材料却很少保存那些选择在城邑层面如何被执行。
\end{causalchain}

统一之后，秦没有停在战场上。度量衡、文字、道路、郡县、法律和巡行被迅速推向新得到的土地。里耶秦简中的县级文书，让后人看见这一速度并非传说：原楚地已经被写进秦的行政日常。帝国的厉害之处，不只在它能打到哪里，而在它能多快让远方开始使用同一套文书、名目和命令。

\begin{sourcenote}[统一不是单一人物的一步棋]
\sourcebook{史记·秦始皇本纪}与《六国年表》给出韩、赵、魏、楚、燕、齐覆亡的编年骨架，《韩世家》《赵世家》《魏世家》《楚世家》《燕召公世家》《田敬仲完世家》以及《白起王翦列传》补入各国视角；《资治通鉴》可供核对其后编年的串联。这些皆为后出的传世叙事，不是六场战役留下的连续军政档案。王翦“六十万”、大梁受灌的水道与时长、燕辽东残部的行军路线、齐廷人物的劝说和每国降附的城数，都须保留为有待检验的细节；现代研究通常结合战场地理、水利条件、交通与粮运，以及战国国家动员能力来衡量它们。里耶秦简、巡游刻石和度量衡器物能够旁证统一后行政与标准化的实际推进，却不能倒过来证实灭国战争中的每一个数字和场景。
\end{sourcenote}

\phantomsection\label{fig:vol01:warring-qin:qin-unification}
\begin{center}
\begin{tikzpicture}[x=.82cm,y=.82cm,font=\small]
  \fill[paper] (0,0) rectangle (12,6.3);
  \fill[dynasty!26] (.8,2.0) -- (3.7,1.8) -- (4.0,4.8) -- (1.4,5.1) -- cycle;
  \node[align=center] at (2.3,3.4) {秦\\关中};
  \fill[yellow!24] (4.3,3.4) -- (5.4,3.3) -- (5.6,4.5) -- (4.5,4.7) -- cycle;
  \node at (4.95,4.0) {魏};
  \fill[green!20] (4.7,4.7) -- (6.8,4.6) -- (7.0,6.0) -- (5.1,6.1) -- cycle;
  \node at (5.9,5.35) {赵};
  \fill[orange!20] (7.1,4.3) -- (8.8,4.3) -- (9.0,5.8) -- (7.3,6.0) -- cycle;
  \node at (8.0,5.05) {燕};
  \fill[cyan!20] (9.2,3.2) -- (11.3,3.3) -- (11.4,5.2) -- (9.4,5.4) -- cycle;
  \node at (10.3,4.3) {齐};
  \fill[red!17] (5.4,.7) -- (9.4,.8) -- (9.6,3.4) -- (6.3,3.6) -- cycle;
  \node at (7.4,2.0) {楚};
  \fill[purple!17] (5.5,2.9) -- (6.5,2.9) -- (6.7,3.9) -- (5.7,4.1) -- cycle;
  \node at (6.1,3.45) {韩};

  \draw[-{Stealth[length=2mm]},ultra thick,color=dynasty] (3.4,3.7) -- (5.4,3.6);
  \draw[-{Stealth[length=2mm]},ultra thick,color=dynasty] (4.4,3.0) -- (5.7,3.35);
  \draw[-{Stealth[length=2mm]},ultra thick,color=dynasty] (5.9,3.1) -- (7.2,2.55);
  \draw[-{Stealth[length=2mm]},ultra thick,color=dynasty] (6.2,4.2) -- (7.2,4.85);
  \draw[-{Stealth[length=2mm]},ultra thick,color=dynasty] (8.5,4.45) -- (9.1,4.35);
  \node[text=dynasty,below] at (4.3,3.1) {秦逐次东进};
\end{tikzpicture}

\smallskip
\small\textit{图  秦统一六国空间示意：展示关中秦与韩、赵、魏、楚、燕、齐的相对位置及东进方向；不表示逐年战线、精确道路或疆界。}
\end{center}


\subsection{帝国的日常，先于帝国的崩溃}

秦统一后最容易被记住的是严刑、徭役和焚书；但如果只写这些，便无法解释它为何能在短时间内统治广阔地域。战场上的降书要变成仓中的出入、县廷的案件和一段可以通行的道路，才算真正进入帝国。秦制所改变的，正是这些日常：户籍怎样被登记，仓储怎样被核验，文书怎样传递，地方官怎样回报，以及度量衡怎样让相隔千里的交换与征发可以互相折算。

齐降后的廷议首先在重新规定最高权力。朝臣讨论“王”“帝”乃至“泰皇”的称谓，王政最终定为始皇帝，\eventanchor{WQ-0221-IMPERIAL-TITLE-DEBATE}{秦·始皇帝称号之议}随后以“朕”“制”“诏”重定命令语言，\eventanchor{QIN-0221-IMPERIAL-DOCUMENTS}{秦·皇帝文书}并以玺印和诏书格式辨别授权的真伪与等级，\eventanchor{QIN-0221-IMPERIAL-SEALS}{秦·印玺与诏书}。这些安排不能被误解为一夜间遍及所有乡里的新语法；它们首先让从中心发出的命令有了可识别的形式。廷议中关于封建功臣还是置郡县的对立，\eventanchor{WQ-0221-FENGJIAN-JUNXIAN-DEBATE}{秦·封建郡县之议}同样由后出的史家压缩而成，但不复封诸侯的制度选择以及中央文武事务的分工，\eventanchor{QIN-0221-THREE-DUKES-NINE-MINISTERS}{秦·中央官署分工}确实构成了新政权运转的框架。

这个框架必须先处理可能重新成为政治中心的人和物。传统的三十六郡数目概括了首轮郡级编制，\eventanchor{QIN-0221-THIRTY-SIX-COMMANDERIES}{秦·三十六郡}却不是一张前221年即永久不变的地图；重大申诉和刑狱须经廷、御史与郡县的文书层级上达，\eventanchor{QIN-0221-LEGAL-APPEALS}{秦·郡县申诉}地方间仍有时限与传递的摩擦。秦又迁徙六国豪富、王族和部分旧贵族，\eventanchor{QIN-0221-RELOCATION}{秦·徙豪富}\eventanchor{WQ-0221-SIX-KING-RELOCATION}{秦·迁六国旧族}，处理旧国宗庙与陵墓，\eventanchor{QIN-0221-OLD-KING-TOMBS}{秦·重组旧国礼仪}并收缴、集中六国兵器，\eventanchor{QIN-0221-WEAPONS}{秦·收缴六国兵器}\eventanchor{WQ-0221-ARSENAL-COLLECTION}{秦·咸阳兵器集中}，又传有拆除内地旧国险塞、强调边塞防务的政策，\eventanchor{QIN-0221-WALL-DEMOCRACY}{秦·内地拆塞}。十二万户、销兵铸像、拆墙范围和具体处置方式主要来自传世记载，不能拿来计算全国人口、军械或工程量；其共同指向是削弱旧国在原地重组兵员、门客和礼仪的资源。

统一并不只靠禁绝。半两钱成为法定钱币形制，\eventanchor{QIN-0221-HALF-LIANG}{秦·半两钱}秦又以律令、权量和文书格式重新接合旧国行政接口，\eventanchor{WQ-0221-LEGAL-UNIFICATION}{秦·律令与标准统一}在主要道路布置传舍、亭和邮人，\eventanchor{QIN-0221-ROAD-POST-STATIONS}{秦·道路传舍}\eventanchor{WQ-0221-COUNTY-POSTAL-NETWORK}{秦·郡县邮传}。这些制度要由县廷实际办理。迁陵县前221年至前208年的简牍档案，\eventanchor{QIN-0221-LIYE}{秦·里耶县廷文书}正留下收发、仓储、邮传和司法事务的局部痕迹：它能证实行政被做成了手续，不能把一个洞庭郡县直接扩写为全国的执行图。陵园中持续生产的陶俑和兵器，\eventanchor{QIN-0221-TERRACOTTA}{秦·兵马俑作坊}也显示大型工程的分工与标准化，不能据此推定所有官营作坊均以同一方式运行。

前220年前后，帝国更具体地把人、道路与都城接在一起。编户、里伍和郡县文书重新登记旧国民众，\eventanchor{QIN-0220-IMPERIAL-CENSUS}{秦·编户登记}但里耶、睡虎地都是局部材料，不能倒推出一次全国同步的普查。皇帝新名分通过诏令、刻石和官府文书反复出现，\eventanchor{QIN-0220-INSCRIPTION-ORDERS}{秦·刻石诏令}御史、使者则巡察新郡县的守令、仓库和法律执行，\eventanchor{QIN-0220-JUNXIAN-INSPECTORS}{秦·巡察郡县}。驰道沿线须勘定路径、处理水沟桥梁和征地，\eventanchor{QIN-0220-ROAD-SURVEYORS}{秦·驰道勘路}又持续征发劳役，\eventanchor{QIN-0220-ROAD-CORVEE}{秦·驰道劳役}；咸阳吸纳豪富、工官和官署，\eventanchor{QIN-0220-XIANYANG-URBAN-EXPANSION}{秦·咸阳扩张}使都城成为枢纽，也更依赖周边的水、粮与人力。

巡行使这些线路变得可见。前220年至前219年，始皇巡经陇西、北地和东方郡县，\eventanchor{QIN-0220-WESTERN-TOUR}{秦·巡行西北与东方}齐鲁守令须为车驾供应食物、警卫和道路，\eventanchor{QIN-0219-JIMO-TOUR-ORDER}{秦·齐鲁巡行供给}。泰山封禅、立刻石，\eventanchor{QIN-0219-TAISHAN-INSCRIPTION}{秦·泰山刻石}琅邪台刻石并留置官吏、迁民，\eventanchor{QIN-0219-LANGYA-INSCRIPTION}{秦·琅邪刻石}\eventanchor{QIN-0219-LANGYA-SETTLERS}{秦·琅邪留民}都是现存或传世的政治文本，却首先是自我宣示。次年张良等在博浪沙袭击车驾未成，\eventanchor{QIN-0218-BOLANGSHA}{秦·博浪沙遇袭}秦在韩、河南旧地搜捕，\eventanchor{QIN-0218-BOLANG-SEARCH}{秦·博浪沙搜捕}并调整车驾与卫队路线，\eventanchor{QIN-0218-IMPERIAL-GUARD-ROUTES}{秦·巡行警卫}说明旧国记忆仍能转化为刺杀和治安风险。

国家的格式最终仍要由会读写、会算期限的人执行。睡虎地墓所见吏卒喜的律令、案例和日书，提示南郡基层吏员如何学习办案，\eventanchor{QIN-0217-NANJUN-LAW-TRAINING}{秦·南郡吏员训练}墓葬于前217年封闭，\eventanchor{QIN-0217-SHUIHUDI}{秦·睡虎地秦墓}其中的律历、日书与期限计算又让抽象法令落到具体日期，\eventanchor{QIN-0217-SHUIHUDI-TERM-CALCULATION}{秦·基层期限计算}。简册的抄写年代并不完全等于墓葬年代，墓主的经历也不等于全南郡；它们仍比宏大诏令更清楚地提醒我们，帝国不是自行运作的法条，而是一连串会误期、会核验、也会被人处理的文本劳动。

齐王建降后，秦王政先以\eventanchor{QIN-0221-EMPEROR}{秦·定号皇帝}为天下重新命名，随即面对一个比称号更棘手的问题：新得的六国土地，应交给旧王族分守，还是交给可以替换、可以考核的官吏？《史记·秦始皇本纪》保存了廷议封建与郡县的著名争论，李斯的发言又把分封说成重启争战的隐患。传世文字未必是廷议的逐句记录，却保留了当时必须作出的制度选择。秦采用\eventanchor{QIN-0221-COMMANDERIES}{秦·推行郡县}：郡统县，中央任命的守、尉、监进入地方，旧国的政治中心不再以世袭封君的名义续存。史书所谓三十六郡，是理解这一转折的入口，不是一张在前221年便永久固定的行政地图；郡数、边界和新占地区的处置仍在调整，县以下的实际运行更不可能由咸阳一纸诏令同时完成。

这项选择的难处，不在于秦廷能否说出“郡县”二字，而在于原有六国的城邑、旧贵族、吏员和赋役关系如何被重新编排。任命守、尉、监，使地方权力在名义上可以由中央替换；它并不自动供给熟悉地方田亩、户口、仓储与道路的人。新制度须借助原有城邑中的人力和材料，又必须让他们使用新的官名、期限与责任链。故而郡县是战后接收的框架，不是接收已经完成的证明；三十六郡的整齐说法尤其不能遮蔽此后对边界、郡数和县级事务的持续调整。

同一时期的\eventanchor{QIN-0221-STANDARDS}{秦·统一文字与尺度}，也不只是后来概括为“书同文、车同轨”的整齐口号。权量器上的诏书铭刻和《秦始皇本纪》所述法度、衡石、丈尺，使一斗粟、一段布、一份仓储簿册有了可由不同官署互认的名目；书写规范则使命令更容易越过旧国的文字习惯。它们服务于市场，也服务于赋役、军需和司法中的计算。可是，标准的颁布不等于地方习惯当日消失，器物、书体和执行速度在各地并不齐一；现存权量与传世文字能证实国家意图和部分实践，不能替我们画出一幅毫无摩擦的“统一瞬间”。

标准化的行政意义在于使不同地方的记录可以被比较和追问。仓中报出的粟、工役所交的物料、县廷判定的价值，若没有相互承认的尺度和书写格式，便很难沿着层级上行；若能互认，中央才可能把远方的一项短缺转换成可核验的数字与责任。权量铭文为这种意图留下了较近时的物证，但它们多是器物和诏令的见证，不是每个市、每个仓的执行报告。所谓“统一”因此应理解为反复推行、校验与纠偏的过程，而非一项诏书抵达便完成的事实。

前220年修治\eventanchor{QIN-0220-ROADS}{秦·驰道}，让这种互认的尺度有了移动的身体。《秦始皇本纪》以“堑山堙谷”的简短笔法写道路工程，后世又将驰道想象为笔直铺向四方的帝国筋骨；哪些路段何时完成、路宽和役夫数目如何，材料远没有这幅图像清楚。可以把握的是，路不再只为一支东进军队而修：使者携文书，地方输送粮械，皇帝与随从巡行，北方驻军等待补给，都要依赖一连串可被维护、被查验的通道。战争中把人粮推向前线的本领，开始被改作使县廷按期把消息和物资送回中心的本领。

驰道也不能被理解为一条从咸阳直达天下尽头、随时畅通的单线。它应当连同既有道路、渡口、山口、驿传和沿线城邑来看；山谷需要施工，河流与季节会限制运输，维护本身还要继续征用劳力与物料。传世材料用宏大的工程动词概括这些困难，后出的道路叙事又常把路线画得过于确定。较稳妥的说法是，秦确在统一后推进以中心为指向的交通干线，借此增加巡行、传令和转运能力；具体哪一段在何时贯通、宽度及工役人数，则仍须等待考古、地理与文献互证。

这种转换并不意味着战争已经退场。前219年，\eventanchor{QIN-0219-TOUR}{秦·东巡刻石}把皇帝带到泰山、琅邪等旧国空间；巡行须由沿途郡县安排道路、供给和警卫，刻石则把“天下一家”的语言留在可见之处。石刻是接近当时的政治文本，却首先是自我宣示，其传世文字也经过抄录与校勘，不能据以断定每一处地方都已心悦诚服。随后北方的\eventref{QIN-0215-XIONGNU}{秦·北逐匈奴}和南方的\eventref{QIN-0214-SOUTH}{秦·经略岭南}又把道路、军粮与郡县延展到更远的边缘。统一后的行政日常，因而从一开始就和边疆战争缠在一起。

\subsection{道路走到边境，簿籍落在县廷}

前215年，蒙恬北逐匈奴并经营河套，\eventanchor{QIN-0215-XIONGNU}{秦·北逐匈奴}。这不是把一条北方边线在地图上画直即可完成的事。河套及其通道的控制，要靠驻军、粮运、城塞与郡县之间不断的往返；从内地运出的军粮若不能持续抵达，军事推进便会变成孤军停驻。长城和北部郡县的连接工程因而既是防务设施，也是转运与征发系统的一部分。传世史书所给的战果与工程叙事可以说明秦将统一后的资源投入北方，却不能让我们精确复原每一段墙体、道路的修筑先后，或把后世所见全部遗存都归到同一年的同一计划。

南方的限制不同。前214年前后，秦军向岭南推进，设置南海、桂林、象等郡，\eventanchor{QIN-0214-SOUTH}{秦·经略岭南}。山地、水系和湿热环境使这里不能照搬关中至中原的调度节奏；军队与地方官能否取得粮食、传递讯息并在雨季维持道路，都是疆域扩展的实际条件。有关灵渠的开凿年代、作用范围以及各郡的具体建置，文献与研究仍有争议，故不能把一项水利工程写成南征所有困难的万能答案。能够确认的是，南扩把郡县和军役推到新的环境，也把长期驻军、运输和与当地群体互动的成本带进帝国。

边疆的军需最终仍会回到县廷的簿籍。睡虎地秦简所见律令、问答和办理案件的文书范式，能看见官吏如何以罪名、期限、验核和连带责任来分类人事；墓葬在统一后不久封闭，简册本身却有战国末至帝国初年的不同层次，不能把其中每一条都当作前221年刚刚颁行的全国法典。里耶秦简则多为前222至前208年迁陵县的实际文书。这个原属楚地的县，已在处理收发、簿籍、调发和上行下达；一根根简牍不讲统一的英雄故事，却让人看见旧楚居民、县吏和远方中央怎样被同一套期限与格式拴在一起。

所谓登记，在此不该被写成人人在同一日被帝国看见的全景图。睡虎地、里耶保存的是有限地点与有限事项中的名籍、收发、检验和责任记录；它们足以说明行政依赖把人、物和期限写入可追查的格式，却不能为全国每户复原一份完整档案。诏令规定什么与基层实际做到什么之间，隔着抄写、送达、误期、核验和回报的每一道工序。恰是这些不显眼的工序，让“郡县”“驰道”与“北逐、南进”不止是朝廷名词，也让一处延误、匿报或人力不足可能沿着同一网络放大。

迁陵县留下的简牍尤其使这种放大变得可见。现存材料中有收取、发遣、转送和答复的痕迹：一份文书并非抵达县廷便告结束，它还须被识别、登记、交给下一位经手者，并在期限内成为可向上说明的事项。正因为它们是零散而具体的办理记录，里耶简比宏大的帝国叙事更能提示郡县如何落在桌案上；也正因为它们只来自一个原楚地县和少数保存下来的年份，不能据此替任何一郡虚构一张完整的日常流程表。睡虎地简所见律令与里耶简所见实务相互照亮，却不是同一时间、同一地点、同一层级的两份全国总档。

睡虎地和里耶都只是保存下来的局部，不足以替全国每一郡县作证，也不能把某份残简的程序直接扩展为所有人的生活。不过它们改变了读《史记》的角度：法律不是只在刑场显形的威吓，文书也不是从中央单向落下的命令。它们要求基层不断登记、封检、转送和答复；只有这些反复的劳动能够接住上面的尺度、道路和巡行。帝国的能力由此不再只是攻城时的兵员数，而是使远方一个县将事情写成中心能够读懂、能够追问的格式。

前216年令黔首自实田，\eventanchor{QIN-0216-LAND-SELF-REPORT}{秦·自实田}可见国家希望取得赋役与土地纠纷的登记口径；此令是否创设私田、地方如何申报，史家和研究者并无定论。同一时期咸阳的兵器集中之后还须维修、再铸和分配，\eventanchor{QIN-0216-URBAN-ARMAMENT-CONTROL}{秦·咸阳军械管理}不能把“收兵”想成一项无须工官、材料和边军补给的象征行动。

北逐之后，河套经营更不可能只留下一个战报。秦迁徙谪戍与民众充实北方郡县，\eventanchor{QIN-0215-FRONTIER-RESETTLEMENT}{秦·北方迁戍}军队又依靠渡口、粮仓和后方道路维持推进，\eventanchor{QIN-0215-HETAO-GRAIN-ROUTE}{秦·河套粮道}迁民与戍者的种子、粮食和工具同样需要长期供应，\eventanchor{QIN-0215-HETAO-MIGRATION-LOGISTICS}{秦·河套迁民供给}。始皇东巡至碣石立刻石，\eventanchor{QIN-0215-JIESHI-INSCRIPTION}{秦·碣石刻石}秦则连接旧燕赵塞、修补或新筑若干段落，\eventanchor{QIN-0215-NORTHERN-WALL-SECTIONS}{秦·北方塞防}并在接管的边地调整郡县、戍卒和既有边防经验，\eventanchor{QIN-0215-YAN-ZHAO-FRONTIER-INTEGRATION}{秦·整合燕赵边地}。这些措施的年代、墙体范围和迁徙户数都不可精确合算，但共同说明边界须以居民、转运和日常维护来保存。

岭南的接管又是另一种缓慢。桂林、象郡在前214年前后列入郡级名目，\eventanchor{QIN-0214-GUILIN-XIANG-COMMAND}{秦·置桂林象郡}秦郡县与当地部众围绕道路、粮食和劳役发生持续协商与冲突，\eventanchor{QIN-0214-LINGNAN-LOCAL-NEGOTIATION}{秦·岭南地方协商}谪戍者、军人和部分编户迁入，\eventanchor{QIN-0214-LINGNAN-MIGRANT-FAMILIES}{秦·岭南迁民}，并非一纸郡名即可代替日常供给。灵渠联结湘、漓水系的传统年代可作为后勤线索，\eventanchor{QIN-0214-LINGQU}{秦·灵渠}其水运仍须配合陆路转运和闸坝维护，\eventanchor{QIN-0214-LINGQU-PORTAGE}{秦·灵渠转运}不能消除山地、雨季和地方抵抗的成本。

南海郡尤其依赖珠江水系与海岸航路，\eventanchor{QIN-0214-NANHAI-SEA-ROUTES}{秦·南海水路}番禺等城邑作为军政与海陆转输的节点，\eventanchor{QIN-0214-PANYU-ADMIN}{秦·番禺郡治}赵佗受命为南海尉，\eventanchor{QIN-0214-ZHAOTUO}{秦·赵佗守南海}。早期航线、郡治和赵佗职掌都未留下可供精确复原的连续档案；秦亡后南越能够自立，仍提示边郡经营积累了不完全由咸阳支配的人员与资源。

这也是它最危险的地方。战国时代，国家把社会资源推到战争前线；统一以后，战争的技术并未被解除，而是进入工程、边防、巡行和日常考核。到\eventanchor{QIN-0213-EMPIRE}{秦·高强度治理}所处的前213年前后，朝廷又以\eventanchor{QIN-0213-BOOKS}{秦·禁书令}压缩民间师承和议论空间，并以\eventanchor{QIN-0212-PALACE}{秦·大工程与坑儒叙事}留下更沉重的征发记忆。淳于越主张师古封建、李斯以郡县反驳的廷议，\eventanchor{QIN-0213-CHUNYUYUE-DEBATE}{秦·封建郡县再议}以及由此进入高层叙事的法令与地方批评，\eventanchor{QIN-0213-QIN-LEGAL-CRITICISM}{秦·法令批评}都经史家压缩；焚书令在郡县的查验和执行也不可能证明为全国同步，\eventanchor{QIN-0213-BOOK-BAN-LOCAL-ENFORCEMENT}{秦·禁书令的地方执行}。禁书政策确有其事；“坑儒”的对象、人数和经过则主要经由汉代叙事被塑形，不能轻率化作系统清除现代意义儒者的确证。

扶苏因谏止严厉处置而被遣往上郡监蒙恬军，\eventanchor{QIN-0212-FUSU-NORTH}{秦·扶苏北监}他与咸阳之间因此更受使者、距离和文书控制的限制，\eventanchor{QIN-0212-NORTHERN-HEIR-ISOLATION}{秦·继承人远离中枢}却不能说已经完全断讯。与此同时，咸阳宫苑、阿房前殿与陵园竞争工匠、木材和徭役，\eventanchor{QIN-0212-XIANYANG-PALACE-LABOR}{秦·咸阳营造劳役}新迁人口与工官集中又使水、粮和劳力调度紧张，\eventanchor{QIN-0212-XIANYANG-POPULATION-CONTROL}{秦·咸阳供给压力}。蒙恬主持由九原通咸阳的直道，\eventanchor{QIN-0212-ZHIDAO}{秦·修筑直道}沿线还须不断维修、供给和征役，\eventanchor{QIN-0212-ZHIDAO-MAINTENANCE}{秦·直道维护}；它缩短军讯距离，也把负担推向沿线社会。

前211年东郡陨石刻有反秦文字的故事，\eventanchor{QIN-0211-METEOR}{秦·东郡陨石谣言}文字本身无从核验，秦廷迁徙或惩治周边居民的说法亦不能化为可计量的司法事实，\eventanchor{QIN-0211-METEOR-COLLECTIVE-PUNISHMENT}{秦·东郡集体惩治}却足以显示异常、流言和皇帝合法性在高压政治中如何被连在一起。道路、工程、巡行、军役和法律也不是一张可以互换的暴政清单，而是同一台机器在不肯降速时，连续落到无数家庭的时间、物资和行动上的要求。

禁书令尤其需要区分命令与执行。传世叙事保存了朝廷限制私藏《诗》《书》及诸子文本、以法令为师的主张，能说明政治批评与统一法令之间的紧张；它不能证明典籍在一夜间焚尽，也不能让后世据此断定所有地方的师承同时中断。宫殿、陵墓与道路工程也是同理：朝廷确能集中劳力和物资，考古与文献却未提供一份可以逐日相加的全国工役总账。把所有压力归为一项工程或一个禁令，会丢掉更关键的制度事实——同一批基层户、吏与转运网络，可能在不同季节被要求供役、输送、戍守和承担法律责任。

\begin{debate}[诏令、传说与执行]
《史记》让禁书、坑儒、阿房宫和骊山工程在同一段落中构成秦亡前的强烈图景；汉代政治论说又乐于将其收束为暴政的证据。这个图景保存了高强度动员和思想控制的历史问题，却不等于一份同期执行簿。禁书的范围与地方落实、阿房宫工程的实际进度、坑儒的对象及人数，都须分别以传世叙事、考古和简牍材料校读。就现有材料而言，压力的存在较为清楚，压力在何处、何时和以何种强度落下，则不能伪造为精确统计。
\end{debate}

\begin{sourcenote}[从皇帝诏令到县廷竹简]
\sourcebook{史记·秦始皇本纪}提供定号、郡县、尺度、驰道和巡行的叙事骨架，《李斯列传》保存郡县廷议的后出传记版本；它们可以显示朝廷如何解释自己的选择，不能复原每次廷议的原话或全国建置的逐月过程。秦权量铭文与泰山、琅邪等刻石使统一标准和巡行的政治语言有了较近时的物证，但铭文的流传、文本校勘及其覆盖范围仍须分别判断。睡虎地秦简所见律令和文书范式、里耶秦简所见迁陵县前222至前208年的收发与簿籍，则把视野压低到地方操作；它们尤其适于说明法律与行政如何运行，不适于据少数保存地点推定全国执行完全一致。现代关于郡县建置、驰道路向、秦律编纂和简牍年代的研究，正是在这些传世叙事、器物和局部档案之间校读，而非用其中一种材料独断全貌。
\end{sourcenote}

这几类材料并不在回答同一个问题。《史记》是西汉史家对统一与覆亡的回望，最适合搭起行动者、事件次序和朝廷自我解释的骨架，也最须防备把后出的廷议、人物言辞当作现场记录。睡虎地、里耶的秦代文书离办理程序更近，却只记录局部事务；权量器、刻石以及宫殿、城塞等考古遗存能检验标准、工程和物质投入留下的痕迹，却通常不能告诉我们某名官吏在某日为何作出选择。至于贾谊《过秦论》以及汉代不断重述的禁书、坑儒与秦亡故事，则是新王朝借秦讨论统治的后设叙事。把它们并置，不是把四种证词揉成一种“秦的真相”，而是让各自能够证明的范围互相限制。

\subsection{巡行的终点与继承的断裂}

\begin{event}{公元前210年}{沙丘继承：巡行队伍里的新皇帝}{可靠纪年，文书细节有争议}{东巡归途与沙丘一带}\eventanchor{QIN-0210-SUCCESSION}{秦·沙丘继承}
秦始皇死在巡行归途中，帝国最重要的命令中心一时不在咸阳，而在一支必须继续移动、又必须控制消息的队伍里。《史记·秦始皇本纪》《李斯列传》《赵高列传》叙述赵高、李斯与胡亥改变继承安排，扶苏被迫自杀，胡亥即位为二世。年代主线和继承结果较为稳固；遗诏如何传递、各人说过什么、谁先得知死讯等细节，主要出自汉代写成的本纪与列传，并无现存秦廷档案可以逐句核对。

最后一次东巡本身已经拉长了中枢的指挥线：车驾经旧齐、旧楚地域，郡县要供应食物、警卫和道路，\eventanchor{QIN-0210-EASTERN-TOUR-ADMIN}{秦·末次东巡供给}会稽刻石则把皇帝礼制投向东南，\eventanchor{QIN-0210-KUAIJI-INSCRIPTION}{秦·会稽刻石}。路线与每日行止并不完整，刻文也有传本问题；可见的是移动的皇帝必须越来越依赖一条由地方接力维持的命令链。

始皇死于沙丘后，谁掌握玺印、诏书与死讯，谁便能先定义合法命令，\eventanchor{QIN-0210-SHAQIU-SEAL}{秦·沙丘掌玺}赵高、李斯又在秘不发丧的情形下控制车驾西返，\eventanchor{QIN-0210-SHAQIU-COURT-MOVEMENT}{秦·沙丘秘行}为改写继承留出时间。鱼鲍掩臭之类细节显然经过文学化；但行宫、卫队和文书在有限人群中被控制，正是这场政变能够发生的物质条件。

扶苏收到矫诏时难以独立核验父皇死讯与中央命令，\eventanchor{QIN-0210-FUSU-EDICT-VERIFICATION}{秦·扶苏难验诏书}胡亥继位后又拘系蒙恬，\eventanchor{QIN-0210-MENGTIAN-EXEC}{秦·扶苏死与蒙恬失势}并逐步剥夺蒙毅等先帝近臣解释遗命的资格，\eventanchor{QIN-0210-MENGYI-REMOVAL}{秦·蒙氏近臣被逐}。扶苏的反应、蒙恬劝谏与诏书形成过程主要见于传记，不宜把它们补成实时对话；其后果却很清楚：北方军政与咸阳之间的一组制衡力量被拆开。胡亥返咸阳以二世名义接收宗庙、官署和宫禁，\eventanchor{QIN-0210-HUHAI-ACCESSION-RITUAL}{秦·胡亥即位}形式上的继承完成了，筛选奏报的权力也更集中到近臣手中。

这场继承危机的约束首先是空间性的。皇帝猝死时，巡行所携带的符玺、随从、文书和交通资源集中在有限的人群手中，北方的扶苏与各地郡县只能通过被截断或被重写的命令得知结果。二世即位并不意味着郡县机器立刻停转；相反，它仍能征发、审讯、派将并催促工程。但信息进入中心的路径越来越窄，能校正决策的人越来越少。对一套依赖文书期限和责任追究的制度而言，这种断裂尤其危险：地方仍收到命令，却较难确信中心是否掌握前线、粮运和民情的真实情形。

不必把秦亡归结为一封被篡改的诏书。沙丘继承是使后续危机加速的触发因素，能够使它扩大为帝国失控的条件，则是前述边防、工程、征发与关东接管所累积的负担仍在运行。李斯、赵高和胡亥在传世叙事中有鲜明而对立的形象；从可核实的后果看，更重要的是扶苏被排除、二世即位后中央的信任结构迅速恶化，面对多地压力时，既有的命令链不再可靠。
\end{event}

\subsection{大泽以后}

\begin{event}{公元前209年}{大泽乡的雨与第一面反旗}{可靠纪年，细节有争议}{蕲县大泽乡一带}\eventanchor{QIN-0209-FALL}{秦·起事与覆亡}\eventanchor{QIN-0209-CHENSHENG}{秦·大泽起事}
一队被征调北上戍守的徒众困在大泽乡的雨水和泥路之间。传统叙事把他们的领头人叫作陈胜、吴广：行期被阻，畏惧军法，因而先杀领队、举兵反秦。这是《史记·陈涉世家》最有力量的开端；它使后来每一次读到秦亡的人，都仿佛先看见一条无法准时通过的道路。可是，故事也过于整齐。我们不能由此断言每一个迟到者必定立即处死，或把传记中激昂的言辞当作大泽乡当日逐字可闻的口号。睡虎地、里耶所见的秦律与县廷文书，能够说明期限、徭役、连带责任和文书考核确实构成严密的日常压力，却不能替这场暴雨逐项作证。

雨阻失期而起事的缘由，\eventanchor{QIN-0209-RAIN-DELAY-CRISIS}{秦·大泽失期危机}正要放在这个证据限度内理解：地点、人数以及“失期即死”是否为当时成文法，均有争论；被征者在期限与灾害之间缺少回旋余地的结构性压力则不可忽略。起事者攻取城邑时杀逐守令，\eventanchor{QIN-0209-CHENSHENG-LOCAL-MAGISTRATES}{秦·郡县守令失守}使吏民不得不选边，原先依赖守令和乡里协作的县制节点遂开始连锁断裂。

陈胜在陈地建立张楚名号，\eventanchor{QIN-0209-ZHANGCHU}{秦·张楚建号}试图把戍卒抗命转成能够吸纳旧国民众的政治语言；随即分遣周文、武臣、周市等多路进军，\eventanchor{QIN-0209-CHENSHENG-COMMANDS}{秦·张楚分兵}又须为扩张中的队伍征粮、募卒，\eventanchor{QIN-0209-ZHANGCHU-TAX-SUPPLY}{秦·张楚军需}。任命的先后、各军路线和财政措施只保留在零散叙事中，不能重建为一份完整军政图；快速分兵超过陈地仓储、文书和统帅能力，却足以解释各路军头为何很快取得自主性。

反秦的组织也从来不只出于陈地。刘邦在沛起兵，\eventanchor{QIN-0209-LIUBANG}{秦·刘邦起兵}萧何、曹参等熟悉县廷的人掌握户籍、仓储和乡里联络，\eventanchor{QIN-0209-LIUBANG-PEI-RECORDS}{秦·沛县吏员网络}使一支地方队伍可以持续募兵、征粮。项梁在会稽起兵，\eventanchor{QIN-0209-XIANGLIANG}{秦·项梁起兵}接收郡守的兵器、仓储与吏民网络，\eventanchor{QIN-0209-XIANGLIANG-KUAJI-ARMORY}{秦·项氏接收会稽}。这些开端的具体计谋带有后出英雄叙事的加工，但它们共同显示，秦的县郡网络一旦失灵，也能为不同集团提供可夺取的粮仓、文书和人手。

对咸阳而言，周文西进至戏地附近，\eventanchor{QIN-0209-ZHOUWEN}{秦·周文逼近关中}使关中东门突然受威胁；二世朝廷从工程、刑徒与关中人力中加派军役，\eventanchor{QIN-0209-QIN-RESPONSE-LEVIES}{秦·应急征发}暂时增加兵源，也进一步加重社会负担。危机同时改变中枢的人事关系：赵高借战局罗织李斯罪状，\eventanchor{QIN-0209-LISI-IMPEACHMENT}{秦·赵高构陷李斯}终于取代李斯控制政务、任为丞相，\eventanchor{QIN-0209-ZHAOGAO-XIANG}{秦·赵高为相}。诬告书、人物动机和任官程序不可尽信于传记；不过一面动员地方应急、一面缩窄能够校正军情的中枢，正使帝国越来越难将命令转成有效回应。

即使将戏剧化的细部暂且搁置，起事何以能在数月之间越出一支小队，仍须解释。触发因素是征戍途中失去退路；使能条件则是统一战争和工程动员留下的役夫、兵器、道路与聚集的人口，和前210年\eventref{QIN-0210-SUCCESSION}{秦·沙丘继承}之后中央命令链日益不可信的局面。陈胜、吴广借“楚”的名号与扶苏、项燕等仍被记得的名字召集人众，并非旧楚自动复活，而是新帝国辖下的人能够借旧国记忆给彼此一个共同的旗帜。郡县的仓廪、亭舍和道路原为征发而设；一旦基层官吏、豪强、亡命者和被征者各自改换计算，它们也使消息与队伍比咸阳的命令传得更快。

这正是秦末最初的命令崩塌：并非所有县都同时倒向反秦，而是中央再不能确信一纸文书抵达后仍会按原意执行。陈胜在陈地称王，吴广的队伍旋即发生内讧，显示新起的政治体同样没有稳定的军令和供给；但失败没有把已经打开的通路关上。沛地的刘邦、楚地项梁及其侄项羽，以及各处以旧国名义集合的人，先后成为独立的行动者。秦面对的已不再是一案可以审结的逃亡，而是许多地方武装在同一交通网中争夺城邑、粮食和名分。
\end{event}

\begin{event}{公元前208年}{章邯出关，陈胜败亡}{可靠纪年}{关中、陈地与河北方向}\eventanchor{QIN-0208-ZHANGHAN}{秦·章邯出关}
二世朝廷并非毫无反应。章邯奉命整编骊山刑徒及其他可用人力出关，先后击败陈胜集团，再向赵地推进。把这支军队简单叫作“囚徒军”会把它写得太像一支临时乌合之众；传世史书所说的刑徒、奴产子与征发者，较可能指向由工程劳动力、受刑者和临时动员者拼接而成的战时兵源，具体编制却难以复原。秦仍能从关中调出人、粮与将领，正说明帝国的动员能力没有在大泽一夜消失。

周文逼近关中后，章邯将骊山刑徒、工程劳力等编为应急军，\eventanchor{QIN-0208-ZHANGHAN-EMERGENCY-ARMY}{秦·章邯应急军}击败周文，\eventanchor{QIN-0208-ZHANGHAN-ZHOUWEN}{秦·章邯破周文}暂时守住关中东门。兵员构成和人数都带有叙事夸张，仍可见常备军分散于边地时，中央只能将本来用于工程和惩罚体系的人力改作野战军。章邯随后攻陈，张楚失去最初的都城与仓储中心，\eventanchor{QIN-0208-ZHANGHAN-CHEN-SIEGE}{秦·章邯攻陈}陈胜在败走中为庄贾所杀，\eventanchor{QIN-0208-CHENSHENG-DEATH}{秦·陈胜败亡}；末路细节有异文，但陈胜之死并未把已经分散到各地的反秦战争收回。

反秦者自身也无从自动结成一国。陈胜对周文、武臣等将领难以维持奖惩、粮饷与信息控制，\eventanchor{QIN-0208-CHENSHENG-COURT-FAILURE}{秦·张楚中枢失灵}秦军反攻时各路力量早已各自为政。章邯在定陶击杀项梁，\eventanchor{QIN-0208-XIANGLIANG-DINGTAO}{秦·定陶杀项梁}楚军遂由项羽、刘邦等分头行动，熊心的名义位置上升；战役过程虽有不同叙述，项梁之死迫使楚地重组力量的后果较为清楚。

章邯再围赵王歇于巨鹿，\eventanchor{QIN-0208-JULU-SIEGE}{秦·围巨鹿}本欲先摧毁赵地的反秦中心，反而迫使诸侯考虑救赵，形成次年会战的条件。与此同时，赵高以丞相身分限制二世接见群臣，\eventanchor{QIN-0208-ZHAOGAO-CHANCELLOR}{秦·赵高隔绝群臣}关东的军情、财政与赦令更难在中枢获得交叉校验。宫禁程序无法逐日复原；帝国在外线不断增兵、在内廷不断收窄信息的两种趋势，却已彼此叠加。

但这也是它的代价正在回流的时刻。常规地方行政不足以处理多点起事，朝廷不得不将修筑、戍守和惩罚体系中的人力改作野战军；地方上每一次镇压又会制造新的离散者和新的武装。陈胜败死，吴广早已不在，反秦者却没有随之归零。赵地被围、楚军重组，项羽在叔父项梁死后取得更重的军事位置；刘邦则在沛地队伍与关中道路之间寻找可进入帝国腹心的路线。人事的重心已从咸阳的诏令，移到谁能替一支队伍供粮、守城并让盟友承认其名义。
\end{event}

\begin{event}{公元前207年}{巨鹿、降军与咸阳的失语}{可靠纪年}{巨鹿、咸阳}\eventanchor{QIN-0207-JULU}{秦·巨鹿之战}
章邯在赵地围困的巨鹿，原本可能成为秦以一场决定性会战重新接管关东的机会。项羽率楚军北上救赵，诸侯军亦在周边观望；《史记·项羽本纪》以“破釜沉舟”写楚军的决绝，留下极鲜明的英雄画面。较可把握的主线是：赵地围困迫使分散的反秦力量集中，秦军久战的补给和内部协调却愈发困难，项羽的胜利因而改变了联军的力量秩序。至于每一支诸侯军何日畏缩、楚军毁掉多少器具，仍是汉代叙事塑造项羽的笔法，不宜拿来填满战场的空白。

楚军能够北进，先有项羽杀宋义、取得救赵主力指挥权，\eventanchor{QIN-0207-XIANGYU-SONGYI}{楚·项羽夺军}后有传世叙事所谓破釜沉舟、强迫军队决战，\eventanchor{QIN-0207-POFU}{楚·破釜沉舟}。宋义与项羽的冲突动机、楚军毁弃器具的程度，都不能从后出的英雄叙事中直接量化；较可靠的是楚军不再观望而投入巨鹿，从而改变了秦军原本分割反秦力量的计算。

巨鹿之后，章邯与王离所部的局面已难维持。章邯终于降项羽，\eventref{QIN-0207-JULU}{秦·巨鹿之战}所暴露的不只是一次战败，更是秦最能野战的军队脱离了咸阳的指挥。关中朝廷此时又在自行拆毁信息与信任。赵高借审讯和控告除去李斯，\eventanchor{QIN-0207-LISI}{秦·李斯被杀}；《史记·李斯列传》《赵高列传》把酷刑、供词和宫廷对答写得极其具体，基本事实可以显示相国失势、赵高专权，细节却带有后出传记将失败人格化的浓重痕迹。对仍在前线等待粮饷与指示的人来说，更要紧的是：负责整合行政与决策的中心，再也无法稳定地协调战争。

章邯在巨鹿败后与项羽议降，须同时权衡楚赵压力和二世朝廷的猜忌，\eventanchor{QIN-0207-QIN-ZHANGHAN-NEGOTIATION}{秦·章邯议降}其书信、地点和心理动机多为传记所塑；秦关东主力正式在洹水一带降楚的结果则较为明确，\eventanchor{QIN-0207-ZHANGHAN-SURRENDER}{秦·章邯降楚}。二世在关东战败、李斯被杀后愈少直接接见将相，赵高垄断奏报，\eventanchor{QIN-0207-QIN-ER-SHI-COURT-ISOLATION}{秦·二世隔绝将相}于是前线、财赋与赦令之间失去稳定的校正渠道。

赵高继而杀二世，立子婴为“秦王”而不复称帝；降格的名号本身承认了天下已不再听命于一个皇帝。子婴不久杀赵高，\eventanchor{QIN-0207-ZIYING}{秦·子婴即位}，使宫廷暂时摆脱了控制者，却不能召回已降的主力，也不能让函谷关以东重新服从。这里没有一条由一名奸臣直通秦亡的因果线：赵高的操弄是加速中央失语的因素，能够使它致命的，是继承危机、战争消耗、地方倒戈和指挥链断裂已经互相叠加。

赵高杀二世而立子婴为秦王的改号，\eventanchor{QIN-0207-ZIYING-KINGSHIP}{秦·子婴称秦王}未能缩小已成事实的危机。刘邦则由峣关、武关方向推进：他在峣关以谈判、奇兵等方式瓦解守军，\eventanchor{QIN-0207-LIUBANG-YAOGUAN}{汉·取峣关}经间道与招降越过关中东南门，\eventanchor{QIN-0207-LIUBANG-XIAO-ROUTE}{汉·峣武间道}\eventanchor{QIN-0207-LIUBANG-WUGUAN-NEGOTIATION}{汉·武关招降}并破武关入关，\eventanchor{QIN-0207-LIUBANG-WUGUAN}{汉·破武关}。张良计谋和守将反应有显著的传记加工，仍不能抹去南部关隘已同中枢失联、难获外援的局面。

子婴仓促在蓝田集结残军阻截，防线终告瓦解，\eventanchor{QIN-0207-LANTIAN-DEFENSE}{秦·蓝田防线瓦解}使刘邦得以接近咸阳。战斗地点、将领与先后在史籍中并不一致；不过秦军主力既降、南门又开，关中已没有一条可以独自支撑帝国的外层防线。

\end{event}

\begin{event}{公元前206年}{刘邦先入关，子婴降秦}{可靠纪年}{武关、霸上、咸阳}\eventanchor{QIN-0206-SURRENDER}{秦·子婴降汉}
项羽在河北承接诸侯军的拥戴时，刘邦所率的队伍由南路逼近关中。两人的力量并不相同：项羽有巨鹿后形成的盟主声望和较大的联军，刘邦的优势在于较早进入秦的南面门户，并能沿途争取或迫使守军让路。后世常把这段竞逐写成两位未来对手预先布好的棋局；在当时，他们面对的更直接问题是道路、关隘、降者和粮食，谁先抵达咸阳，谁便能取得处置秦廷和关中资源的发言权。

刘邦至霸上，子婴以秦王身分出降，帝国至此终结。投降仪式的细节、刘邦入咸阳后的禁令和项羽随后入关的冲突，主要由《史记·高祖本纪》与《项羽本纪》保存；二书让后来楚汉相争的胜者与败者从一开始便有鲜明的性格轮廓，必须防止把汉初以后形成的对照直接投回前206年。仍可确定的是，子婴所能交出的已是一座失去关东主力和外部命令网络的都城；秦的终结不是某一座宫门被攻破的瞬间，而是数年间地方不再执行、野战军离散、朝廷相互猜忌之后，关中终于没有足以拒绝一支先到军队的防线。

刘邦入关后的首要问题是让城市与粮食市场不再随占领而崩散。他封府库、约束军纪以安定关中供给，\eventanchor{QIN-0206-GUANZHONG-FOOD-SECURITY}{汉·安定关中粮市}又以约法三章作出不同于秦严法的政治姿态，\eventanchor{QIN-0206-YUEFA}{汉·约法三章}。禁掠的实际范围与民众反应无法量化，萧何收取丞相府律令、户籍和地图档案的叙事亦有传记塑形，\eventanchor{QIN-0206-LIUBANG-QIN-ARCHIVES}{汉·接收秦档案}却提示占领者知道人口、道路和仓储信息与兵器同样重要。

项羽处置秦降军和旧都的方式构成了另一种压力。新安处置章邯旧部，传称坑杀二十余万降卒，\eventanchor{QIN-0206-XINAN-PRISONERS}{楚·新安降卒}人数与过程不能当作统计，秦人对楚军的敌意因大批降卒的命运而加深则合乎战后处境。项羽入咸阳后杀子婴、焚掠宫室的记载，\eventanchor{QIN-0206-XIANGYU-XIANYANG}{楚·项羽入咸阳}同样应区分较可靠的入关与子婴被杀，和受汉代史家影响的焚掠规模、责任叙事；秦地官民随后面对的占领与惩罚秩序，\eventanchor{QIN-0206-XIANGYU-GUANZHONG-PUNISHMENT}{楚·处置秦旧都}却是无可回避的现实。

鸿门宴中的项、刘谈判没有消除关中处置权的冲突，\eventanchor{QIN-0206-HONGMEN}{楚·鸿门谈判}其对白和座次虽经文学化，双方力量并存的危机并非虚构。项羽把关中分为雍、塞、翟三国，\eventanchor{QIN-0206-THREE-QINS-ENFEOFFMENT}{楚·分封三秦}再把刘邦封为汉王、迫其赴汉中，\eventanchor{QIN-0206-LIUBANG-HANZHONG-TRANSFER}{汉·刘邦赴汉中}企图安置秦降将并隔开新的竞争者；分封边界、兵员和路线有异说，重造地方军事中心的方向则十分清楚。

项羽入关后的分封、焚掠与诸军处置立刻把“反秦”联盟变成新的权力竞逐；那已是楚汉战争的开端。秦灭亡的直接后果，却并非郡县、法律、度量衡和文书从此消失。刘邦和后来建立的汉朝必须接收关中仓储、道路、官吏经验与统治尺度，也必须回答一个更难的问题：怎样继承\eventref{QIN-0213-EMPIRE}{秦·高强度治理}所证明的国家能力，又不把同样密集的征发和惩罚重新压回社会。

因此，县吏、律令和仓储技术被新统治者选择性保留和改写，\eventanchor{QIN-0206-QIN-LEGAL-CONTINUITY}{汉·选择性继承秦制}“约法三章”并不能证明秦律一概废除。秦的王朝覆亡了，能够让赋役、治安与军粮重新运转的行政信息却成了继承者不得不使用、也不得不重新限制的制度资源。
\end{event}

\begin{sourcenote}[谁把秦亡写成了这个样子]
\sourcebook{史记·陈涉世家}、\sourcebook{史记·项羽本纪}、\sourcebook{史记·高祖本纪}与\sourcebook{史记·秦始皇本纪}提供了从大泽到子婴降秦的基本叙事骨架，\sourcebook{资治通鉴}将其编入逐年脉络；它们都不是秦末现场留下的连续军政档案。《史记》写成于汉武帝时，陈胜的豪言、项羽的勇决、刘邦的宽厚或机变，以及赵高的阴险，都服务于人物列传和王朝兴亡的解释。秦简能校正我们对法律和地方行政的想象，却多属不同地区、不同时间的碎片，不能直接证明大泽每一项情节。因而，陈胜吴广首先举事、章邯出军、巨鹿后秦主力崩解、刘邦入关与子婴投降的年代主线较为稳固；口号、对白、人数和单一动机，则应保留为汉代回望中的叙事，而非未经媒介的秦末声音。
\end{sourcenote}

\begin{textwindow}[贾谊《过秦论》]
“仁义不施而攻守之势异也。”

出处为贾谊的《过秦论》。这篇汉初政论回望秦亡，意在让汉廷避免重蹈覆辙；它把秦的兴亡压缩成一句有力判断，不能替代前文所见的继承危机、战争消耗和地方倒戈，却能让读者看见后人怎样以秦的命运讨论新帝国的统治。
\end{textwindow}

\begin{turningpoint}[制度得失：秦留下的不是一个答案]
秦用战国形成的动员能力完成统一，解决了长期分裂、标准不一和诸侯战争的问题；代价是中心把统一后的社会仍按高强度战时逻辑对待。秦亡证明这种逻辑难以长期承受，却没有证明统一帝国本身不能存在。汉以后所有王朝都要在秦留下的能力与代价之间重新作答。
\end{turningpoint}
